\documentclass[10pt,a4paper]{amsart}
\usepackage[margin=19mm]{geometry}
\usepackage{amsmath,amssymb,mathtools}
\usepackage[scr=rsfs,cal=euler]{mathalfa}
\usepackage{xfrac}
\usepackage[unicode,hidelinks,bookmarks=false]{hyperref}
\newcommand{\ie}{i.e.\ }
\newcommand{\cf}{cf.\ }
\newcommand{\resp}{respectively\ }
\newcommand{\Subsection}{\subsection}
\providecommand{\nobreakdash}{\nobreak\hbox{-}}
\setlength{\emergencystretch}{2em}
\multlinegap0pt
\usepackage{mathtools}
\usepackage{amssymb}
\usepackage{bm}
\usepackage{bbm}
\newtheorem{lemma}{Lemma}
\newcommand{\Om}{\Omega}
\newcommand{\kappam}{\kappa_m}
\title{\Large \textbf{Feedback Stabilization and Tracking for Heat Equations}\\ \textbf{Using Thermo-Plasmonic Nanoparticles as Actuators}}
\author{Arpan Mukherjee, S\'ergio S. Rodrigues, Mourad Sini}
\date{Source: arXiv:2602.14581v1 (CC BY 4.0)}
\begin{document}
\maketitle

\noindent\emph{Source and licence.} \Large \textbf{Feedback Stabilization and Tracking for Heat Equations}\\ \textbf{Using Thermo-Plasmonic Nanoparticles as Actuators}. Version arXiv:2602.14581v1 (CC BY 4.0), distributed by arXiv under the Creative Commons Attribution 4.0 International licence, http://creativecommons.org/licenses/by/4.0/, as recorded in the arXiv licence metadata for this exact version and shown on the abstract page https://arxiv.org/abs/2602.14581v1. Full licence notice: \texttt{licenses/arxiv-2602.14581-CC-BY-4.0.txt}.\par\smallskip

\noindent\textit{Editorial scope.} Counted unit: the article's lemma lem:tail-small-compact (source0.tex lines 877--911). The article's complete original proof of it is delivered with it (source0.tex lines 913--971, one \texttt{proof} environment of 59 lines). No mathematics is written, repaired or re-derived: the statement and the proof are the article's own lines, in the article's own notation, and no retained line is substituted or re-flowed.  Retained uncounted as the source-local context the proof needs: lines 656--661; lines 783--787; lines 829--860.  Nothing was omitted from any retained span, so the package carries no omission record.  No retained line contains a citation, so the wrapper carries no bibliography and no bibliography entry.  No retained line contains a figure environment, an image-inclusion command or drawing code, so no image is delivered and the package declares no asset dependency.  Editorial formatting only, in addition: an AMS \texttt{amsart} preamble that restates the article's own declarations and macros used by the retained lines, namely the environments \texttt{lemma} on separate counters, together with the standard packages those lines need.  The retained environments print in this document as Lemma 1 in order of appearance, whereas the article prints the same items with the numbering of its own full document.  The complete native source of the article as distributed in the arXiv e-print for this exact version is bundled unchanged as \texttt{sources/194722/source0.tex}.
\begin{equation}
\label{eq:3.2}
H := L^2(\Om),\qquad
A_0 y := \kappam\Delta y,\quad
D(A_0) := \{ y\in H^2(\Om): \partial_\nu y = 0 \text{ on }\partial\Om\}.
\end{equation}

\begin{equation}
\label{eq:def-C-full}
C:H\to\mathbb{R}^M,\qquad
(Cz)_j := \big(\mathcal A^{-1}z\big)(x_j),\quad j=1,\dots,M.
\end{equation}

\begin{lemma}[Matrix representation of $C$ and low/high decomposition]
\label{lem:C-matrix}
Let $z=\sum_{k\ge1}\alpha_k\varphi_k\in H$. Then
\[
\mathcal A^{-1}z = \sum_{k\ge1}\frac{\alpha_k}{1+\lambda_k}\,\varphi_k,
\]
and, for each $j=1,\dots,M$,
\[
(Cz)_j = \sum_{k\ge1} d_{jk}\,\alpha_k,\qquad
d_{jk} := \frac{1}{1+\lambda_k}\,\varphi_k(x_j).
\]
Define the (infinite) matrix $D=(d_{jk})_{1\le j\le M,\ k\ge1}$. Then
\[
Cz = D\alpha,\qquad \alpha:=(\alpha_1,\alpha_2,\dots)^\top\in\ell^2.
\]
For the fixed $M$ (number of actuators and low modes), decompose
$D=(D_{MM}\;|\;D_{M>})$ with
\[
D_{MM} := (d_{jk})_{j,k=1}^M\in\mathbb{R}^{M\times M},\qquad
D_{M>} := (d_{jk})_{1\le j\le M,\ k>M}.
\]
Writing $\alpha=(\alpha_M,\alpha_>)$ with
\[
\alpha_M := (\alpha_1,\dots,\alpha_M)^\top,\qquad
\alpha_> := (\alpha_{M+1},\alpha_{M+2},\dots)^\top,
\]
we obtain
\begin{equation}
\label{eq:Cz-split}
Cz = D_{MM}\alpha_M + D_{M>}\alpha_>.
\end{equation}
\end{lemma}

\begin{lemma}[Compactness of the measurement operator and small tail]
\label{lem:tail-small-compact}
Let $\Omega\subset\mathbb{R}^3$ be a bounded $C^2$ domain, $A_0$ the Neumann
Laplacian as in \eqref{eq:3.2}, and $\{\varphi_k\}_{k\ge1}$ an
$L^2(\Omega)$–orthonormal basis of eigenfunctions of $-A_0$ with eigenvalues
$0=\lambda_1\le\lambda_2\le\cdots$. Let $x_1,\dots,x_M\in\Omega$ be fixed
distinct points, and define $C$ as in \eqref{eq:def-C-full}.
\newline
For $\alpha=(\alpha_k)_{k\ge1}\in\ell^2$, set $z=\sum_{k\ge1}\alpha_k\varphi_k$
and write $Cz=D\alpha$ with $D=(d_{jk})_{j,k}$ as in
Lemma~\ref{lem:C-matrix}. For each $N\in\mathbb{N}$, decompose
$D=(D_{MN}\mid D_{M> N})$ where
\[
D_{MN} := (d_{jk})_{1\le j\le M,\ 1\le k\le N}\in\mathbb{R}^{M\times N},\quad
D_{M> N} := (d_{jk})_{1\le j\le M,\ k>N}.
\]
\noindent
Then the operator $T:\ell^2\to\mathbb{R}^M$ defined by $T\alpha:=D\alpha$ is
compact. In particular,
\begin{equation}
\label{eq:Dtail-norm->0}
\|D_{M> N}\|_{\mathcal{L}(\ell^2,\mathbb{R}^M)} \xrightarrow[N\to\infty]{} 0.
\end{equation}
Equivalently, for every $\varepsilon>0$ there exists $N_\varepsilon\in\mathbb{N}$
such that for all $N\ge N_\varepsilon$ and all
$\alpha=(\alpha_k)_{k\ge1}\in\ell^2$,
\[
\big\|D_{M> N}\alpha_{>N}\big\|_{\mathbb{R}^M}
= \big\|C\Big( \sum_{k>N}\alpha_k\varphi_k \Big)\big\|_{\mathbb{R}^M}
\le \varepsilon\,\|\alpha\|_{\ell^2},
\]
where $\alpha_{>N} := (\alpha_{N+1},\alpha_{N+2},\dots)$.
In particular, by choosing the truncation index $N$ sufficiently large, the
measurement of the high-frequency tail can be made arbitrarily small.
\end{lemma}

\begin{proof}
Define
\[
J:\ell^2 \to H,\qquad J\alpha := \sum_{k\ge1}\alpha_k\varphi_k.
\]
Since $\{\varphi_k\}$ is an orthonormal basis of $H=L^2(\Omega)$, $J$ is an
isometric isomorphism from $\ell^2$ onto $H$.
\noindent
Let $R:=\mathcal A^{-1}:H\to D(A_0)\subset H^2(\Omega)$. Since
$H^2(\Omega)$ embeds continuously and compactly into $C(\overline\Omega)$,
then the composite operator
\[
R_J := R\circ J:\ell^2\to C(\overline\Omega)
\]
is compact. The evaluation map
\[
E: C(\overline\Omega)\to\mathbb{R}^M,\qquad
(Ew)_j := w(x_j),\quad j=1,\dots,M,
\]
is bounded. Hence the measurement operator
\[
T := E\circ R_J:\ell^2\to\mathbb{R}^M
\]
is compact as a composition of a compact and a bounded operator.
\newline
By construction, for $\alpha\in\ell^2$ and
$z=J\alpha=\sum_{k\ge1}\alpha_k\varphi_k$,
\[
T\alpha = E(RJ\alpha) = C z = D\alpha,
\]
so $T$ is represented by the matrix $D$ in the basis $\{\varphi_k\}$.
\newline
Let $P_N:\ell^2\to\ell^2$ be the orthogonal projection onto the first $N$
coordinates:
\[
P_N\alpha := (\alpha_1,\dots,\alpha_N,0,0,\dots).
\]
Set $T_N := T P_N$, which has finite rank. Since $T$ is compact, we have
\[
\|T - T_N\|_{\mathcal{L}(\ell^2,\mathbb{R}^M)}
\xrightarrow[N\to\infty]{} 0.
\]
On the other hand,
\[
(T-T_N)\alpha = T(I-P_N)\alpha,
\]
and $(I-P_N)\alpha$ has coordinates
$(0,\dots,0,\alpha_{N+1},\alpha_{N+2},\dots)$, so
\[
(T-T_N)\alpha = D_{M> N}\,\alpha_{>N}.
\]
Hence
\[
\|T-T_N\|_{\mathcal{L}(\ell^2,\mathbb{R}^M)}
= \|D_{M> N}\|_{\mathcal{L}(\ell^2,\mathbb{R}^M)},
\]
which proves \eqref{eq:Dtail-norm->0}. The last statement is just the
operator norm inequality written explicitly.
\end{proof}

\end{document}
